\documentclass[11pt,a4paper]{article}
\usepackage{amsmath,amssymb,amsthm,bm}
\usepackage{geometry}
\geometry{margin=1in}
\usepackage{hyperref}
\usepackage{graphicx}
\usepackage{setspace}
\usepackage{booktabs}
\usepackage{listings}
\onehalfspacing

\title{\textbf{Paper OOL-9: Protocell Integration, Replication, and Error Thresholds:\\ The First Heritable Container Gate}}
\author{Steve Bico Mujjabi, MD\\
\small Independent Researcher\\
\small Founder, Vura Labs\\
\small Kampala, Uganda\\
\small ORCID: \href{https://orcid.org/0009-0001-0556-5516}{0009-0001-0556-5516}}
\date{May 2026}

\begin{document}
\maketitle

\begin{abstract}
This paper is the integration gate. A protocell claim requires bounded copying error, retained internal chemistry, coupled energy throughput, and a boundary whose persistence covaries with its internal state. The paper combines Eigen--Orgel error-threshold logic with a conservative protocell diagnostic and explicitly reports numerical fragility where it appears. That restraint is intentional: Paper 9 defines what must be true before Paper 10 can introduce parasitic exploitation.
\end{abstract}

\section*{Symbol Conventions}

\noindent\textbf{CDFD origins-of-life symbol convention.}
\begin{center}
{\small
\begin{tabular}{p{0.15\linewidth}p{0.34\linewidth}p{0.41\linewidth}}
\toprule
Symbol & Meaning in Part II & Origins-of-life reading \\
\midrule
$\Phi$ & Available driving flux & Redox, thermal, photochemical, proton/ion, or chemical potential drive supplied by the environment \\
$C$ & Effective constraint or resistance & Mineral geometry, pore confinement, diffusion barrier, permeability, activation barrier, or maintenance burden \\
$S$ & Responsive organization & Surface, boundary, catalytic network, coacervate, or protocellular structure that routes flux into retained function \\
$M_s$ & Retained structural memory & Composition, sequence, topology, compartment state, mineral imprint, or boundary history that persists across renewal cycles \\
$\Psi_s$ & Local adaptive ratio $(\Phi/C)\cdot S\cdot M_s$ & Bookkeeping gauge for whether drive, constraint, responsiveness, and retention are co-present; not a universal constant \\
$\Lambda_{\mathrm{life}}$ & Integrated Life Number where used & Capstone surplus ratio for decay, near-critical proto-biology, and sustained self-maintenance \\
\bottomrule
\end{tabular}
}
\end{center}

\noindent\textbf{Greek-symbol convention.}
\begin{center}
{\small
\begin{tabular}{p{0.15\linewidth}p{0.34\linewidth}p{0.41\linewidth}}
\toprule
Symbol & Local role & Interpretation rule \\
\midrule
$\alpha$ & Accumulation or reinforcement coefficient & Defines how fast a constraint, pathway, or retained structure grows under load \\
$\beta$ & Relaxation, decay, or washout coefficient & Defines recovery, loss, clearance, hydrolysis, or return toward baseline \\
$\gamma$ & Coupling, diffusion, or redistribution coefficient & Defines spread across pores, surfaces, compartments, networks, or chemical phases \\
$\delta,\epsilon$ & Perturbation, tolerance, or small offset scales & Used only as local thresholds, numerical guards, or experimental tolerances \\
$\eta,\lambda,\mu,\nu$ & Efficiency, rate, baseline, or frequency terms & Meaning is set by the equation and subscript where they occur \\
$\rho,\sigma,\tau,\phi$ & Density/correlation, conductance/noise, time-scale, or phase/field terms & Never universal constants unless a paper explicitly defines them that way \\
$\Omega,\Lambda$ & Domain/state space; Life Number where used & $\Omega$ names a modeled state space; $\Lambda$ is reserved for life-number notation only where defined \\
\bottomrule
\end{tabular}
}
\end{center}


\section{Introduction}
Papers 1--8 build the prerequisite ladder: thermodynamic framing, geochemical energy and transport, oscillating polymerization, autocatalytic closure, and adaptive boundaries. A protocell requires these pieces to become coupled. Replication without a boundary may preserve molecular information but does not necessarily create a lineage. A boundary without inherited internal chemistry can divide physically without biological heredity.

Paper 9 therefore asks two linked questions:
\begin{enumerate}
    \item When does copying preserve information rather than dissolve into error?
    \item When does copying affect the persistence and division of the container that carries it?
\end{enumerate}

\section{Template Growth and Replication Gain}
For a schematic template $T$ and monomer pool $M_i$,
\begin{equation}
    T+M_i\rightarrow TM_i,
    \qquad
    R_{\mathrm{rep}}=k_{\mathrm{rep}}c_Tc_M.
\end{equation}
Transport-limited monomer supply can be represented phenomenologically as
\begin{equation}
    \bm{J}_M=-\frac{1}{C}\nabla c_M.
\end{equation}
A well-mixed caricature gives
\begin{equation}
\frac{dc_T}{dt}
=
k_{\mathrm{rep}}c_Tc_M-\frac{1}{CL^2}c_T
=
c_T\left(k_{\mathrm{rep}}c_M-\frac{1}{CL^2}\right).
\end{equation}
Define the replication gain
\begin{equation}
    \Gamma_{\mathrm{rep}}=k_{\mathrm{rep}}c_MCL^2.
\end{equation}
The condition $\Gamma_{\mathrm{rep}}>1$ is only a linearized growth condition. It is not sufficient for life.

\section{Error Threshold}
If per-step fidelity is $q$ across $n$ additions, the simplest master-sequence survival approximation is
\begin{equation}
    Q=q^n,\qquad \epsilon=1-Q.
\end{equation}
Information persists only if replication gain outruns error catastrophe. A schematic compatibility condition is
\begin{equation}
    \Gamma_{\mathrm{rep}}Q>1,
    \label{eq:error-threshold}
\end{equation}
with the usual warning that real prebiotic copying is not a clean independent-site process~\cite{Orgel1963,Eigen1971}.

The CDFD hypothesis is that fidelity, monomer activation, retention, and damage are coupled to energy throughput. A useful bookkeeping score is
\begin{equation}
\mathcal{H}_{\mathrm{copy}}
=
\Gamma_{\mathrm{rep}}Q
R_{\mathrm{mono}}
R_{\mathrm{ret}}
R_{\mathrm{damage}}^{-1},
\end{equation}
where $R_{\mathrm{mono}}$ is activated-monomer availability, $R_{\mathrm{ret}}$ is retention of templates and products, and $R_{\mathrm{damage}}$ collects hydrolysis, photodamage, side reactions, and strand loss. This is a checklist, not a calibrated law.

\section{Quasispecies Logic Versus Protocell Logic}
The quasispecies threshold and the protocell threshold are related but not identical. The first asks whether a sequence distribution preserves information. The second asks whether a bounded domain inherits a coupled chemistry that affects its growth or persistence.

The protocell condition can be stated as a covariance requirement:
\begin{equation}
\mathrm{Cov}\!\left(
\mathrm{copied\ internal\ state},
\mathrm{container\ persistence}
\right)>0.
\end{equation}
Without this covariance, selection may act on free templates or on empty containers, but not yet on a protocell lineage.

\section{Integrated Protocell Scaffold}
An integrated toy protocell has state variables such as
\begin{equation}
U=\{c_i,\Phi_e,\Phi_H,C_i,\ldots\},
\end{equation}
evolving through transport, reaction, and constraint laws~\cite{SmithMorowitz2016}. A minimal transport skeleton is
\begin{equation}
\frac{\partial c_i}{\partial t}
=
\nabla\cdot\left(\frac{1}{C_i}\nabla c_i\right)
+R_i(c,\Phi_e,\Phi_H),
\end{equation}
with potentials satisfying analogous updates:
\begin{align}
\frac{\partial \Phi_e}{\partial t}
&=
\nabla\cdot\left(\frac{1}{C_e}\nabla\Phi_e\right)+S_e,\\
\frac{\partial \Phi_H}{\partial t}
&=
\nabla\cdot\left(\frac{1}{C_H}\nabla\Phi_H\right)+S_H.
\end{align}
Flux-coupled reactions can be written as
\begin{equation}
\bm{J}_e=-\frac{1}{C_e}\nabla\Phi_e,\qquad
\bm{J}_H=-\frac{1}{C_H}\nabla\Phi_H,
\end{equation}
\begin{equation}
R_i=f_i(c)\,g(|\bm{J}_e\cdot\bm{J}_H|),
\end{equation}
for nonnegative $g$.

\section{Numerical Audit}
The diagnostic for this paper combines a reduced error-threshold sweep with a
bounded two-dimensional protocell update. It should be read as a numerical
stress test, not as validation. The reduced diagnostic reported:
\begin{itemize}
    \item $\Phi=0.2$: final master-sequence fitness $0.367$,
    \item $\Phi=0.5$: final master-sequence fitness $0.083$,
    \item $\Phi=1.0$: final master-sequence fitness $0.014$,
    \item $\Phi=5.0$: final master-sequence fitness $1.000$,
    \item $\Phi=10.0$: final master-sequence fitness $1.000$.
\end{itemize}
The non-monotonic low-to-intermediate response rules out a single universal critical flux in this reduced diagnostic. The correct claim is windowed: copying benefits from energy throughput only when activation, retention, and damage remain bounded.

The final release replaces the earlier unbounded explicit autocatalytic update
with saturated kinetics and bounded adaptive $S,M_s$ updates. Under those
release parameters, the two-dimensional protocell diagnostic remains finite for
the full run. The result should still be read conservatively: the simulated
boundary protects a low internal $\Psi_s$ region while the exterior experiences
large drive, which is a stress-test picture of compartmental buffering rather
than a demonstration of real protocell reproduction.

\begin{figure}[ht]
\centering
\includegraphics[width=\textwidth]{outputs/paper_09/error_threshold_and_protocell.png}
\caption{Merged error-threshold/protocell diagnostic. The left panel tests flux-dependent master-sequence persistence. The right panel shows a bounded two-dimensional protocell stress test with finite internal and external adaptive-ratio traces. The figure is a mechanism diagnostic, not a successful protocell simulation.}
\end{figure}

\section{Hardening Roadmap}
The release implementation now uses saturated kinetics such as
\begin{equation}
R_{\mathrm{auto}}
=
\frac{k\Phi^2}{K+\Phi^2}-\ell\Phi,
\end{equation}
and verifies boundedness in the generated summary JSON. The next hardening step
is stronger: timestep refinement, grid-size refinement, and comparison against
operator-split or implicit updates for coupled transport and reaction. The
model should only be promoted from stress test to prediction if those
refinement checks preserve the qualitative result.

\section{Experimental Targets}
The model protocell literature gives concrete targets: template-directed genetic-polymer synthesis in fatty-acid vesicles~\cite{Mansy2008}, nonenzymatic RNA synthesis inside model protocells~\cite{Adamala2013}, multi-step nonenzymatic primer-extension chemistry~\cite{Walton2018}, and recent water-compatible thioester-mediated RNA aminoacylation/peptidyl-RNA chemistry~\cite{Singh2025}. CDFD should therefore be tested by a coupled score,
\begin{equation}
\mathcal{R}_{\mathrm{proto}}
=
\Gamma_{\mathrm{rep}}Q
R_{\mathrm{retention}}
R_{\mathrm{exchange}}
R_{\mathrm{division}},
\end{equation}
not by replication alone.

\section{Conclusion}
Paper 9 is the integration gate. It establishes that bounded copying, retained
internal chemistry, coupled energy flow, and heritable containers must be tested
together. The current numerical status is deliberately conservative: the
reduced copying toy is informative, while the two-dimensional protocell
implementation is a bounded stress test that still needs refinement before it
can be treated as a prediction. Paper 10 can only discuss parasitic exploitation
after this replication/container machinery exists.

\section{Reproducibility Note}
The companion script \texttt{supplementary\_ool\_09.py} writes the figure
\path{outputs/paper_09/error_threshold_and_protocell.png} plus CSV/JSON
summaries. No notebook is required.
\bibliographystyle{plain}
\bibliography{ool_references}
\end{document}
